\chapter{Gelecek Çalışmalar}
\label{ch:gelecek-calismalar}

Bu bölüm, kapsam dışı bırakılan ama üzerinde düşünülmüş iki konuyu ele
alır: gerçek INT8 nicemleme ve ham görüntülerin bir ağ bağlantısı
üzerinden gerçek zamanlı işlenmesi. İkisi de ``denenmedi'' olarak
işaretlenmiştir -- aşağıdaki hiçbir sayı bir ölçüm değil, bir plandır.

\section{Gerçek INT8 Nicemleme: NVIDIA TensorRT Model Optimizer}
\label{sec:int8-gelecek}

\Cref{sec:hassasiyet}'te sunulan INT8 sonuçları PyTorch içinde bir
\textbf{simülasyondur}: her katmanın çıktısı 8 bitlik bir ızgaraya
yuvarlanıp geri float'a çevrilir, yalnızca etkinleşimler üzerinde --
ağırlıklara dokunulmaz ve hiçbir TensorRT motoru derlenmez. Bu, INT8'in
vereceği zararın bir \textbf{alt sınırını} verir, bir tahmini değil.

\textbf{Planlanan, henüz başlanmadı: NVIDIA TensorRT Model
Optimizer}~\cite{nvidia2024modelopt} (\texttt{nvidia-modelopt}, eski adıyla
AMMO). TensorRT'in doğrudan gerçek INT8 çekirdeklerine derleyebileceği bir
ağ üretmek için NVIDIA'nın PyTorch tarafı PTQ/QAT aracı. Mekanik olarak:

\begin{enumerate}
  \item Her \texttt{nn.Linear} / \texttt{nn.Conv2d} vb.\ katmanı sahte-nicemleme
        Q/DQ düğümleriyle sarmalar -- \textbf{ağırlıklar dahil}, kanal
        başına;
  \item her Q düğümünün ölçeğini, temsili girdi üzerinden kısa bir
        geçişle (maks / entropi / yüzdelik dilim kalibratörleri)
        \textbf{kalibre eder};
  \item kalibre edilmiş modeli Q/DQ düğümleri gömülü halde ONNX'e aktarır;
        TensorRT'in ayrıştırıcısı bu düğümleri doğrudan tanır (``açık
        nicemleme'') ve etraflarına yerel INT8 çekirdekleri derler;
  \item isteğe bağlı QAT: yalnızca PTQ yetersizse Q/DQ etkinken ince ayar
        yapılır.
\end{enumerate}

\textbf{Durum: başlanmadı.} \texttt{pip install nvidia-modelopt}, bir
kalibrasyon veri kümesi, bir Q/DQ ONNX'e aktarım, yeni motor derlemeleri
ve bunlar hazır olduğunda bir \texttt{check\_trt\_parity.py} /
\texttt{compare\_backends.py} koşusu gerektirir.

\section{Ham Görüntülerin Ethernet Üzerinden Gerçek Zamanlı İşlenmesi}
\label{sec:ethernet-gelecek}

\Cref{sec:read-disarida}'da belirtildiği gibi, bu raporda alıntılanan
kare bütçesi \textbf{kare zaten bellekteyken} tracker'ın maliyetidir --
``model yeterince hızlı mı'' sorusunun doğru cevabıdır, ama henüz
\textbf{ölçülmüş bir uçtan uca dağıtım rakamı değildir}. Hedeflenen
kurulumda kareler bir kamera/ağ bağlantısından Orin'e Ethernet ile
gelecektir; bu bağlantı henüz kurulmadığı için hariç tutulan
\textasciitilde35~ms'lik dosya-çözme maliyetinin yerine gerçekte neyin
geçeceği bir ölçüm değil bir tahmindir:

\begin{table}[H]
  \centering
  \caption{Hariç tutulan maliyetin yerine geçmesi beklenen (ölçülmedi).}
  \begin{tabular}{ll}
    \toprule
    Kareler nasıl gelirse & 34~ms'in yerine geçen \\
    \midrule
    Ham, bellekte & $\sim$0,001~ms (arabelleği sarmalama) / 1,4~ms (bir kez kopyalama) \\
    Sıkıştırılmış & bir çözme -- libtiff'ten ucuz ama sıfır değil \\
    \bottomrule
  \end{tabular}
\end{table}

Seçim serbest değildir: 2048×1536, 16 bitte kare başına 6,3~MB'dir; 30
FPS ham veri \textasciitilde1,5~Gbit/s gerektirir -- 1~GbE'nin
taşıyabileceğinden fazla. Ya bağlantı 10~GbE olmalı, ya kareler
sıkıştırılmalı, ya da her karenin daha azı gönderilmelidir (bu,
\cref{sec:gercek-kayitlar}'daki \texttt{crop512} için ikinci bir
gerekçedir: model kadar telin üzerinden geçmesi gerekeni de küçültür).
Sıkıştırma seçilirse, Orin'in donanımsal NVJPG çözücüsü bu maliyeti
CPU'dan alabilir -- ama bu da ölçülmesi gereken bir şeydir, varsayılacak
değil.

\textbf{Planlanan sonraki adım.} Ethernet bağlantısı kurulduğunda,
soketten kareyi alıp modele verene kadar geçen süreyi
\cref{sec:olcum-yontemi}'deki aynı disiplinle (ölçülür, raporlanır,
gerekiyorsa ayrı kategoriye konur) ölçmek, ve bu sayıyı
\cref{sec:uctan-uca-verim}'deki bütçeye ekleyerek gerçek uçtan uca
rakamı elde etmek.

\section{Diğer Açık Sorular}

\begin{itemize}
  \item \textbf{\texttt{--opt-level 5} bellek dikkatine yardımcı olur
        mu?} Bellek dikkati FLOP'ların \%61,9'udur ve dört motor
        arasında en düşük hızlanmaya (1,85×) sahiptir
        (\cref{sec:motor-dogruluk-hiz}). Yalnızca o motoru bu seviyeyle
        yeniden derlemek, TensorRT'in daha agresif bir arama
        yapıp yapmadığını gösterir.
  \item \textbf{Tek motora birleştirme yeniden değerlendirilmeli mi?}
        \Cref{sec:fuzyon-karari}'de ölçülen 0,217~ms'lik başlatma
        boşluğu şu an eşiğin altındadır; bu, mimari değişmediği sürece
        sabit bir sonuçtur, ama farklı bir donanımda ya da motor
        sayısı arttığında yeniden ölçülmeye değer.
  \item \textbf{Gerçek kayıtlarda kırpma hedefi kaybediyor mu?}
        \Cref{sec:gercek-kayitlar}'daki \texttt{crop512} videoları
        izlenerek, IoU'ya değil gözle görülür başarısızlığa
        bakılmalıdır.
\end{itemize}
